\documentclass[10pt,a4paper]{amsart}
\usepackage[margin=19mm]{geometry}
\usepackage{amsmath,amssymb,mathtools}
\usepackage[scr=rsfs,cal=euler]{mathalfa}
\usepackage{xfrac}
\usepackage[unicode,hidelinks,bookmarks=false]{hyperref}
\newcommand{\ie}{i.e.\ }
\newcommand{\cf}{cf.\ }
\newcommand{\resp}{respectively\ }
\newcommand{\Subsection}{\subsection}
\providecommand{\nobreakdash}{\nobreak\hbox{-}}
\setlength{\emergencystretch}{2em}
\multlinegap0pt
\usepackage{bm}
\usepackage{bbm}
\usepackage{dsfont}
\usepackage{xspace}
\usepackage{cancel}
\usepackage{mathtools}
\usepackage[shortlabels]{enumitem}
\usepackage{amsthm}
\makeatletter
\@ifundefined{theorem}{\newtheorem{theorem}{Theorem}}{}
\@ifundefined{lemma}{\newtheorem{lemma}[theorem]{Lemma}}{}
\makeatother
\makeatletter
\newcommand{\cay}[2]{\mathsf{Cay}\left(#1 ; #2\right)}
\expandafter\ifx\csname labelledtheorem\endcsname\relax
\newenvironment{labelledtheorem}[1]
  {\renewcommand\theinnercustomthm{#1}\innercustomthm}
  {\endinnercustomthm}
\fi
\makeatother
\title[Flip colouring of graphs II]{Flip colouring of graphs II}
\author{Xandru Mifsud}
\date{Source: arXiv:2401.02315v3 (CC BY 4.0)}
\begin{document}
\maketitle

\noindent\emph{Source and licence.} Flip colouring of graphs II. Version arXiv:2401.02315v3 (CC BY 4.0), distributed under the Creative Commons Attribution 4.0 International licence, as recorded in the arXiv licence metadata for this exact version and shown on the abstract page https://arxiv.org/abs/2401.02315v3. Full rights notice: licenses/arxiv-2401.02315-CC-BY-4.0.txt.\par\smallskip

\noindent\textit{Editorial scope.} Counted unit: the Lemma whose source label is \texttt{GapsProp} in the article (source0.tex lines 375--379, source label \texttt{GapsProp}), delivered together with the article's own complete original proof of it (source0.tex lines 381--422, one \texttt{proof} environment of 42 lines). No mathematics is written, repaired or re-derived: statement and proof are the article's own text line for line. Required context retained uncounted: SPColourLemma (lines 174--180); CPColourLemma (lines 203--209). Every definition, display and label of those lines is retained unchanged. Omitted from this package and not redistributed: 1--173, 181--202, 210--374, 380--380, 423--476. Nothing inside a retained excerpt is omitted or altered: every retained body is delivered exactly as the article's lines are written, so no omission receipt of any kind accompanies this package. Citations: the retained excerpts carry no citation command at all, so no bibliography entry of the article is transcribed and no reference is redistributed. Reference adaptation: none was needed and none was made; every reference inside the delivered unit resolves inside it, so no unresolved reference or citation is produced. Printed slips kept literal: no printed word, symbol, hypothesis or number of the counted statement or of its proof was altered. Layout and class: the article's own class is \texttt{article} with packages including a4, amssymb, amsmath, amsthm, url, xcolor, enumerate, float, lineno, bezier, amsfonts, graphicx, inputenc, subcaption; the wrapper is the lane's pinned \texttt{amsart} editorial wrapper at 10pt on A4, which restates the theorem-like environments the retained text needs and adds the article's own macro definitions that the retained text needs (\texttt{\textbackslash cay}), with extra wrapper packages (bm, bbm, dsfont, xspace, cancel, mathtools, enumitem). The delivered numbering is the unit's own. Assets: no retained span contains a figure environment, a graphics-inclusion command, a PDF-page inclusion, a table or a TikZ picture; no image file is copied into this package, the package declares no asset dependency and no image file is redistributed with it. Licence: version arXiv:2401.02315v3 of this article is distributed under the Creative Commons Attribution 4.0 International licence, as recorded in the arXiv licence metadata for this exact version and shown on the abstract page https://arxiv.org/abs/2401.02315v3. A full rights notice is delivered at licenses/arxiv-2401.02315-CC-BY-4.0.txt. Characters: every retained excerpt is pure ASCII, so the delivered engine, XeLaTeX with its default Latin Modern text font, reports no missing character.
\begin{lemma}\label{SPColourLemma}
	Let $G$ and $H$ be edge-coloured from $\{1, \dots, k\}$. Then in the coloured strong product $G \boxtimes H$, for any $1 \leq j \leq k$ and $(u, v) \in V(G \boxtimes H)$,
	\begin{enumerate}[i.]
		\item $\deg_j \big((u,v)\big) = \deg_j^H (v) + \deg_j^G (u) \left(1 + \deg^H (v)\right)$,
		\item $e_j \big[(u, v)\big] = e_j^H [v] \left(1 + \deg^G (u)\right) + e_j^G [u] \left(1 + \deg^H (v) + 2\sum\limits_{i = 1}^k e_i^H [v]\right)$.
	\end{enumerate}
\end{lemma}

\begin{lemma}\label{CPColourLemma}
	Let $G$ and $H$ be edge-coloured from $\{1, \dots, k\}$. Then in the coloured Cartesian product $G \ \square \ H$, for any $1 \leq j \leq k$ and $(u, v) \in V(G \ \square \ H)$,
	\begin{enumerate}[i.]
		\item $\deg_j \big((u,v)\big) = \deg_j^G (u) + \deg_j^H (v)$,
		\item $e_j\left[(u,v)\right] = e_j^G[u] + e_j^H[v]$.
	\end{enumerate}
\end{lemma}

\begin{lemma} \label{GapsProp}
	Let $q, k \in \mathbb{N}$ such that $1 < q < \frac{k}{4}$. Let $D_1, \dots, D_q \in \mathbb{N}$ such that $D_q (k - 4q) > 1 + \xi q(q-1) + 5\binom{k-q}{2}$ where $\xi = \max\limits_{1 \leq j < q}\{D_j - D_{j+1}\}$. 
	
	If there exists a $(a_1, \dots, a_q)$-flip graph $F$ such that for every $v \in V(F)$ and $1 \leq j \leq q$, $e_j^F [v] = D_j$, then given any $N \in \mathbb{N}$ there exists a $(a_1, \dots, a_k)$-flip graph for some $a_{q+1}, \dots, a_k \in \mathbb{N}$ where $a_k > N$.
\end{lemma}

\begin{proof}
Let $\Gamma$ be a sufficiently large finite group such that $\Gamma$ has $k - q - 1$ disjoint and inverse-closed subsets $S_j$ where $|S_j| = k - q - j$ for $1 \leq j < k - q$, and such that $S = \bigcup\limits_{j = 1}^{k - q - 1} S_j$ is sum-free in $\Gamma$. 

Consider $K = \cay{\Gamma}{S}$ with an edge-colouring such that the edges labelled in $S_j$ are coloured using $q + j$. Hence for any $v \in \Gamma$, $\deg_{q+j}^K (v) = k - q - j$ and by the sum-free condition on $S$, $e_{q+j}^K [v] = k - q - j$.

Now consider the coloured Cartesian product $F \ \square \ K$, which is $k - 1$ coloured since $F$ is coloured using $1, \dots, q$ and $K$ is coloured using $q + 1, \dots, k - 1$. 

By Lemma \ref{CPColourLemma}, for a given colour $j$ the graph has colour-degree $a_j$ inherited from $F$ for $1 \leq j \leq q$, and $a_j = k - j$ inherited from $K$ when $q < j < k$. Likewise, the number of edges coloured $j$ in a closed neighbourhood is $D_j$ inherited from $F$ for $1 \leq j \leq q$ and $D_j = k - j$ inherited from $K$ for $q < j < k$. 

Finally, note that $F \ \square \ K$ is $\mu = \binom{k-q}{2} + \sum\limits_{i = 1}^q a_i$ regular. Let $t \in \mathbb{N}$ such that $$t \geq \dfrac{1 + \mu + 2\sum\limits_{i = 1}^k D_i}{(k-q)\min\limits_{q < i < k}\{D_i - D_{i+1}\}}$$ and let $H$ be a $\rho = (k-q)t + \binom{k-q}{2}$ regular bipartite graph. For $1 \leq j \leq k - q$, colour $t + j - 1$ matchings of $H$ using $q + j$. 

Let $G$ be the coloured product $H \boxtimes (F \ \square \ K)$. By Lemma \ref{SPColourLemma}, for $v \in V(G)$ and $1 \leq j \leq k$, the edge-colouring in $G$ is such that $$\deg_j (v) = 
\begin{cases} 
a_j & 1 \leq j \leq q \\
a_j + (t + i - q - 1)(1 + \mu)& q < j \leq k
\end{cases}
$$ and $$e_j [v] = \begin{cases} 
D_j (\rho + 1) & 1 \leq j \leq q \\
D_j (\rho + 1) + (t + j - q - 1)\left(1 + \mu + 2\sum\limits_{i=1}^k D_i\right)& q < j \leq k
\end{cases}.
$$

We will show that $G$ as constructed and edge-coloured is a $k$-flip graph. Observe that for $1 \leq j \leq k$, $a_j < 1 + \mu$. Hence for $j = q$ we have that $$\deg_q (v) = a_q < 1 + \mu < a_{q+1} + t(1 + \mu) = \deg_{q+1} (v)$$ and for $j > q$ we have that $\deg_{j+1} (v) - \deg_j (v) = \mu > 0$. Consequently the colour-degree sequence in $G$ is strictly increasing.

Since for $1 \leq j < q$ we have $D_j > D_{j+1}$ then in $G$ we have $e_j [v] > e_{j+1} [v]$. Next note that $\rho + 1 = (k-q)t\kappa$ for some $\kappa > 1$. Hence $(D_q - D_{q + 1})(\rho + 1) > (D_q - D_{q+1})(k-q)t$. Therefore to show that $e_q [v] > e_{q+1} [v]$ it suffices to show $$(D_q - D_{q+1})(k-q) > 1 + \mu + 2\sum\limits_{i=1}^k D_i.$$ 

From the lower bound on $D_q$ in the theorem statement we have that
\begin{align*}
	&D_q (k - q) &&\\
	&> 1 + (3q)D_q + \xi q(q-1) +  5\binom{k-q}{2} &&\\
	&> 1 + \sum\limits_{i = 1}^q a_i + (2q)D_q + \xi q(q-1) +  5\binom{k-q}{2} && \because a_1 < \dots < a_q \leq D_q\\
	&= 1 + \mu + (2q)D_q + \xi q(q-1) +  4\binom{k-q}{2} && \because \mu = \sum\limits_{i = 1}^q a_i + \binom{k-q}{2}\\
	&\geq 1 + \mu + 2\sum_{i=1}^q D_i + 4\binom{k-q}{2} && \because D_{q-i} \leq D_q + i\xi\\
	&= 1 + \mu + 2\sum_{i=1}^k D_i + 2\binom{k-q}{2} && \because \binom{k-q}{2} = \sum\limits_{i = q + 1}^{k} D_i\\
	&= D_{q+1} (k - q) + 1 + \mu + 2\sum_{i=1}^k D_i && \because D_{q+1} = k - q - 1
\end{align*}
as required and therefore the colours $q$ and $q+1$ flip in $G$.

Consider the final case when $q < j < k$. By the choice of $t$ and $\kappa > 1$, we have that $$(D_j - D_{j + 1})(\rho + 1) = (D_j - D_{j+1})(k-q)(t\kappa) > 1 + \mu + 2\sum\limits_{i = 1}^k D_i$$ which we can re-arrange to get $e_j [v] > e_{j+1} [v]$. Hence the sequence of closed neighbourhood sizes is strictly decreasing as required.

It follows that $G$ is a flip graph on $k$ colours, such that for any vertex $v$,  the difference between $\deg_q (v)$ and $\deg_k (v)$ grows in $t$ as $t \to \infty$.
\end{proof}

\end{document}
